\documentclass{situationist-preprint}

\title{Architecture as Public Art}
\msid{AES-17-a}
\msfamily{AES-17}
\msdiscipline{Aesthetics}
\mstype{Principal manuscript}
\mskeywords{architecture; function; dependent beauty; public art; captive audience; context; aesthetic standards}

\begin{document}
\maketitle

\begin{abstract}
Is architecture art? Its buildings serve purposes, are made largely by others than their designers, and are encountered in passing rather than attended to, and these facts have been taken to show that architecture is a craft with aesthetic aspects rather than an art. We argue that architecture is an art, and that its public character imposes stricter standards on it than other arts face. Function does not exclude art: a building's aesthetic properties are those it has as perceived in its category, which includes its function, so that functional beauty is a form of what Kant called dependent beauty, and the design of a building is appreciable as the architect's achievement. Architecture is nonetheless distinctive in having an unavoidable audience. Its works are permanent, occupy shared space and are encountered by people who have not chosen to attend to them. We show that this captive audience removes two licenses other arts enjoy, the license to address only a chosen audience and the license to impose disturbing responses on those who opt in, and adds a requirement of fit with context. A building must therefore meet every standard of art and further ones.
\end{abstract}

\section{The problem}

Vitruvius held that a building must have firmness, commodity and delight, and placed architecture among the liberal arts. Modern aesthetics has been less sure. \citet{Kant1790} classed architecture among the arts of beautiful form, but held that what is essential in it is the suitability of the product for a certain use, so that its beauty is at best dependent beauty, beauty under a concept of what the object is for, rather than the free beauty of a flower or an ornament (\S\S16, 51). Hegel ranked it lowest of the arts, as the art whose material is least able to embody spirit. In the twentieth century the question was posed by the success of functionalism: if form follows function, the designer of a building solves an engineering problem, and whatever beauty results is a by-product.

\citet{Scruton1979} gave the most sustained defence of architecture as an aesthetic practice, while insisting on its differences from the fine arts: it is a vernacular art, practiced by builders as much as by architects, its works are functional, localized and public, and its standards are those of the appropriate rather than of the expressive. \citet{Graham1989} argued that architecture is an art precisely because it is a form of practical problem solving that can be done with more or less imagination. \citet{Goodman1985} argued that buildings mean, by denotation, exemplification and expression, and so are symbolic works in the same sense as paintings. \citet{ParsonsCarlson2008} developed an account of functional beauty on which an object's looking fit for its function is a source of aesthetic value, and \citet{Carlson2000} applied an account of the appreciation of environments to buildings in their settings. \citet{Winters2007} surveys the questions.

We argue for two theses. First, architecture is an art. Second, its public character imposes on it stricter standards than other arts face. The first thesis rejects the view that function excludes art; the second accepts the observation behind it, that architecture is not a fine art like the others, and locates the difference in its audience rather than its function.

\section{Framework}

\begin{definition}[Public character]\label{d:public}
A work has a \emph{public character} if it is (i) permanent on the scale of human lives, (ii) located in space shared by a community, and (iii) encountered, as a condition of ordinary life, by people who have not chosen to attend to it. People who encounter it in this way form its \emph{captive audience}.
\end{definition}

\begin{definition}[Stricter standards]\label{d:strict}
An art $A$ is subject to \emph{stricter standards} than an art $B$ if every condition a work of $B$ must meet to be good of its kind has a counterpart that a work of $A$ must meet, and works of $A$ must meet further conditions, or are denied excuses available to works of $B$.
\end{definition}

\begin{principle}[Art]\label{p:art}
A practice is an art if its products can have aesthetic properties and artistic value fixed by their makers' choices and appreciable as their makers' achievement; and such products, presented within an artworld and taken up by its criticism, are works of art (AES-01).
\end{principle}

\begin{principle}[Category]\label{p:cat}
A work's aesthetic properties are those it has as perceived in its correct category \citep{Walton1970}, and the category of a functional object includes its function \citep{ParsonsCarlson2008}.
\end{principle}

\begin{principle}[Imposition]\label{p:impose}
The license of a work to prescribe a response, including discomfort, revulsion or disquiet, to its audience depends in part on the audience's having chosen to attend to it. Where the audience has not chosen, a work may impose such responses only when they are warranted by the work's public purpose.
\end{principle}

Principle~\ref{p:impose} is a principle about the ethics of address, and it is familiar outside architecture: a loud recording that one may play in one's own room may not be played in a public square, and an image that may be exhibited in a gallery with a warning may not be posted on a billboard. It does not say that disturbing works are bad, which AES-12 denies; it says that their license to disturb has a condition.

\section{The argument}

\begin{lemma}[Function does not exclude art]\label{l:func}
Buildings can have aesthetic properties and artistic value fixed by their designers' choices and appreciable as their designers' achievement.
\end{lemma}

\begin{proof}
By Principle~\ref{p:cat}, a building is perceived in a category that includes its function: as a church, a library, a railway station. Its aesthetic properties as so perceived include the grandeur or intimacy of its spaces, the clarity with which its structure is expressed, the rhythm of its facade and the way its plan articulates the activity it houses. These depend on the designer's choices among the many forms that would serve the function equally well, since function underdetermines form: two libraries of the same capacity and program may be unlike in every aesthetic respect. That function constrains the choice does not make the result less the designer's; sonnets and fugues are made under constraints too, and their constraints are part of what their achievement is measured against. Designing a building whose form both serves and articulates its use, in a situation that could have yielded a merely serviceable one, is an achievement in the sense of AES-04.
\end{proof}

\begin{theorem}[Architecture is art]\label{t:art}
Architecture is an art.
\end{theorem}

\begin{proof}
By Lemma~\ref{l:func} and Principle~\ref{p:art}, architecture is a practice whose products can have the properties required of an art. Its products have been presented and taken up as art by a criticism with its own history, institutions and canon, from Vitruvius to the present, so that the second clause of Principle~\ref{p:art} is met.
\end{proof}

\begin{lemma}[Architecture is public]\label{l:pub}
Architecture has a public character in the sense of Definition~\ref{d:public}, and the other principal arts, in their standard forms, do not.
\end{lemma}

\begin{proof}
Buildings last for generations, stand in streets and squares, and are seen by everyone who passes and used by everyone who works or lives in them, whether or not they have any interest in them as works. A novel is read, a symphony heard and a painting seen by those who choose to read, hear or see them, and each can be avoided at no cost. Public sculpture and murals share the public character, and the conclusions below extend to them; they are exceptions within their arts, while a building's public character belongs to it as a building.
\end{proof}

\begin{theorem}[Stricter standards]\label{t:strict}
The public character of architecture imposes on it stricter standards than the non-public arts face.
\end{theorem}

\begin{proof}
We check the conditions of Definition~\ref{d:strict}. First, every condition a work of a non-public art must meet has an architectural counterpart, since by Theorem~\ref{t:art} buildings are judged as works of art, by their aesthetic merit in their category and their achievement. Second, architecture is denied two excuses. A difficult work in a non-public art may be justified by the rewards it offers to an audience that chooses to make the effort of repeated, informed attention; a building's audience is largely captive (Lemma~\ref{l:pub}) and does not choose that effort, so a building must reward the attention its passers-by and users actually give it, and not only the attention of the qualified observer. And by Principle~\ref{p:impose}, a work that imposes discomfort or revulsion is licensed to do so for a chosen audience, but a building may impose such responses on its captive audience only when its public purpose warrants them. Third, architecture faces a further condition: since a building is permanent and located in shared space, its aesthetic properties include its relation to its surroundings, and it must fit its context, in the sense that it does not destroy aesthetic value already present in the place it occupies without producing compensating value of its own. Non-public works face no such condition, since a painting does not alter the room of anyone who has not chosen to hang it. All three conditions of Definition~\ref{d:strict} are met.
\end{proof}

The theorem explains the distinctive form of architectural controversy. Disputes about buildings are often between those who judge a building as a work, assessed by a qualified observer in the gallery of architectural history, and those who must live beside it. Theorem~\ref{t:strict} says that both kinds of judgment bear on its merit as architecture, and that a building can succeed by the first standard and fail by the second. The controversy over Brutalist housing estates, admired by critics for their structural honesty and resented by many of their residents, is of this kind; the companion manuscript AES-17-b examines it. The account also explains when a building may disturb. A memorial to atrocity, such as Libeskind's Jewish Museum in Berlin, may impose disorientation and unease on its visitors, because its public purpose is the remembrance of what should unsettle; an office block may not, since nothing in its purpose warrants the imposition.

\section{Objections and replies}

\begin{objection}[Dependent beauty is lesser beauty]\label{o:kant}
On Kant's account, the judgment of dependent beauty is not a pure judgment of taste. If architecture's beauty is dependent, architecture is at most an impure art \citep{Kant1790}.
\end{objection}

\begin{reply}
Purity, in Kant's sense, concerns the grounds of a judgment, not the value of its object. Most of the beauty of representational art is dependent in his sense, since it is beauty under a concept of what the work represents, and no one concludes that painting is an impure art. By Principle~\ref{p:cat}, every work is judged in a category, and the category of a building includes its function as the category of a portrait includes its sitter.
\end{reply}

\begin{objection}[Vernacular]\label{o:vern}
Most buildings are designed by no architect, built to pattern and noticed by no one. If architecture is the practice that produces them, it is a craft, and only a few of its products are art \citep{Scruton1979}.
\end{objection}

\begin{reply}
The same holds of other arts: most pictures are diagrams, snapshots and advertisements, and most music is functional. Theorem~\ref{t:art} says that architecture is an art, not that every building is a work of art. The vernacular is relevant in another way, which Scruton was right to stress: it forms the context into which architectural works are inserted, and the condition of fit in Theorem~\ref{t:strict} gives it a claim on them.
\end{reply}

\begin{objection}[Stricter standards stifle innovation]\label{o:innov}
If buildings must please passers-by and fit their contexts, architecture cannot innovate, since innovation is at first unfamiliar and displeasing. The Eiffel Tower was denounced by artists of its day and is now the emblem of its city.
\end{objection}

\begin{reply}
The condition of fit requires that a building not destroy existing value without compensating value; it does not require that it resemble its neighbours, and a building that creates a new focus of value, as the Eiffel Tower did, meets it. The condition on captive audiences requires that a building reward the attention it actually receives over its lifetime, not on the day it opens. Stricter standards are harder to meet, and that is their point; they constrain innovation in architecture as the obligation not to deceive constrains invention in journalism, without excluding it.
\end{reply}

\section{What the result does not establish}

The result does not rank architecture above or below the other arts. It does not say that public standards are set by majority preference: the passer-by's attention is the attention that a building must reward, but whether it does so is judged by the standard of taste of AES-03, applied to that attention. It does not settle which purposes warrant the imposition of disturbing responses beyond the cases discussed, and it does not address the political question of who should decide what is built.

\section{Conclusion}

Architecture is an art. Its function places its works in categories that include their purposes, as representation places paintings in categories that include their subjects, and within those categories its works have aesthetic properties and embody achievements that are their designers'. What distinguishes it is its captive audience. A building is permanent, occupies shared space and is encountered by people who have not chosen to attend to it; it therefore must reward the attention it actually receives, may disturb only when its public purpose warrants it, and must fit its context. Architecture meets every standard of art and further ones.

\begin{thebibliography}{99}\small
\bibitem[Carlson(2000)]{Carlson2000} Carlson, A. (2000). \emph{Aesthetics and the Environment: The Appreciation of Nature, Art and Architecture}. Routledge.
\bibitem[Goodman(1985)]{Goodman1985} Goodman, N. (1985). How buildings mean. \emph{Critical Inquiry}, 11(4), 642--653.
\bibitem[Graham(1989)]{Graham1989} Graham, G. (1989). Art and architecture. \emph{The British Journal of Aesthetics}, 29(3), 248--257.
\bibitem[Kant(1790)]{Kant1790} Kant, I. (1790). \emph{Critique of the Power of Judgment}. Trans. P. Guyer and E. Matthews. Cambridge University Press, 2000.
\bibitem[Parsons and Carlson(2008)]{ParsonsCarlson2008} Parsons, G. and Carlson, A. (2008). \emph{Functional Beauty}. Oxford University Press.
\bibitem[Scruton(1979)]{Scruton1979} Scruton, R. (1979). \emph{The Aesthetics of Architecture}. Princeton University Press.
\bibitem[Walton(1970)]{Walton1970} Walton, K.~L. (1970). Categories of art. \emph{The Philosophical Review}, 79(3), 334--367.
\bibitem[Winters(2007)]{Winters2007} Winters, E. (2007). \emph{Aesthetics and Architecture}. Continuum.
\end{thebibliography}

\end{document}
